% =====================================================================
% tesi/capitoli/30-produrre-dentro-il-gioco.tex
% =====================================================================
% copre: pokedex-home-completo/STUDIO-08-esecuzione-di-codice-e-manipolazione-del-generatore.md
% copre: pokedex-home-completo/STUDIO-08-esecuzione-di-codice-e-manipolazione-del-generatore.md#studio-08-il-terzo-cluster-prioritario-cioè-come-si-produce-dentro-il-gioco-invece-che-intorno
% copre: pokedex-home-completo/STUDIO-08-esecuzione-di-codice-e-manipolazione-del-generatore.md#che-cosa-il-cluster-porta-e-perché-cambia-la-forma-del-lavoro
% copre: pokedex-home-completo/STUDIO-08-esecuzione-di-codice-e-manipolazione-del-generatore.md#la-prima-generazione-e-il-mew-che-ci-manca
% copre: pokedex-home-completo/STUDIO-08-esecuzione-di-codice-e-manipolazione-del-generatore.md#la-seconda-generazione-e-la-via-che-copre-le-nostre-cartucce
% copre: pokedex-home-completo/STUDIO-08-esecuzione-di-codice-e-manipolazione-del-generatore.md#la-terza-generazione-e-la-forma-che-la-via-prende-là
% copre: pokedex-home-completo/STUDIO-08-esecuzione-di-codice-e-manipolazione-del-generatore.md#la-manipolazione-del-generatore-cioè-l-alternativa-legittima
% copre: pokedex-home-completo/STUDIO-08-esecuzione-di-codice-e-manipolazione-del-generatore.md#due-voci-laterali-che-vale-registrare
% copre: pokedex-home-completo/STUDIO-08-esecuzione-di-codice-e-manipolazione-del-generatore.md#che-cosa-ne-discende-per-il-lavoro-e-che-cosa-resta-da-verificare
% copre: pokedex-home-completo/STUDIO-04-la-via-del-dns-e-il-servizio-rianimato.md
% copre: pokedex-home-completo/STUDIO-04-la-via-del-dns-e-il-servizio-rianimato.md#studio-04-la-via-del-nome-di-dominio-e-un-servizio-in-rete-che-qualcuno-ha-rianimato
% copre: pokedex-home-completo/STUDIO-04-la-via-del-dns-e-il-servizio-rianimato.md#che-cosa-è-stato-trovato
% copre: pokedex-home-completo/STUDIO-04-la-via-del-dns-e-il-servizio-rianimato.md#perché-questa-via-non-è-come-quelle-che-stiamo-percorrendo
% copre: pokedex-home-completo/STUDIO-04-la-via-del-dns-e-il-servizio-rianimato.md#che-cosa-è-verificato-e-che-cosa-no
% copre: pokedex-home-completo/STUDIO-04-la-via-del-dns-e-il-servizio-rianimato.md#gli-ostacoli-pratici-che-non-sono-piccoli
% copre: pokedex-home-completo/STUDIO-04-la-via-del-dns-e-il-servizio-rianimato.md#le-due-decisioni-del-progetto-che-questa-via-rimetterebbe-in-discussione
% copre: pokedex-home-completo/STUDIO-04-la-via-del-dns-e-il-servizio-rianimato.md#che-cosa-fare-in-ordine
% verificato-al-commit: ee0f52b
% ---------------------------------------------------------------------

\chapter{Produrre dentro il gioco invece che intorno}
\label{cap:dentro}

I capitoli precedenti hanno prodotto esemplari come file: un programma scrive i byte secondo il modello che il verificatore impiega per giudicare, e la legittimità dell'esito è una proprietà della somiglianza a quel modello. Questo capitolo tratta le due vie che non compongono nulla, cioè quelle in cui l'esemplare è generato dal gioco stesso, e la ragione per cui meritano un capitolo non è la comodità ma la provenienza, che il capitolo~\ref{cap:provenienza} ha stabilito essere una grandezza distinta dai byte.

La distinzione che le governa entrambe si può enunciare in una riga. Un esemplare scritto nella scatola di un salvataggio è \emph{creato}; uno che il gioco consegna perché il suo codice è stato indotto a consegnarlo, o perché una carta vera è arrivata attraverso il canale che il gioco prevede, è \emph{ottenuto}. Le due categorie hanno lo stesso aspetto e storie diverse, e il servizio di destinazione, come si è visto, conosce la storia.

\section{La via interna al gioco, generazione per generazione}

In prima generazione la via è un oggetto malformato ottenibile con una manovra nota, e la procedura documentata la impiega per una cosa sola: riscrivere il nome e l'identificativo dell'allenatore del primo esemplare della squadra portandoli a quelli dell'esemplare distribuito ufficialmente, cosicché il trasferimento lo accetti. La procedura richiede una squadra e uno zaino disposti con quantità precise, perché sono quelle quantità a essere il codice eseguito, e la fonte avverte che l'oggetto impiegato fuori da quegli allestimenti fa cadere il gioco. Una variante fissa i valori individuali in modo da rendere l'esemplare cromatico al trasferimento, con il vincolo che il livello sia almeno cinque.

La lettura di quella voce ha prodotto la scoperta più concreta del suo giro di lavoro, e non era il suo scopo dichiarato. L'esemplare che la procedura imita ha un allenatore e un identificativo precisi, e nelle tabelle del verificatore sta in un file a parte rispetto agli altri doni della prima generazione, come dono unico di livello cinque, con una sola mossa e tutti i valori individuali al massimo. Il lotto prodotto da questo lavoro contiene gli esemplari della medesima specie distribuiti ai tornei e alle manifestazioni, i cui valori individuali sono fissati ad altri numeri, e non contiene questo: è una voce mancante sull'asse delle distribuzioni, ed è la più celebre di tutte. Va registrata anche la conferma incrociata, perché è la prova che la fonte comunitaria è affidabile su questo punto: la guida afferma che i valori individuali sono massimi, e il verificatore lo conferma nel proprio codice assegnando il valore massimo a tutti e cinque quando il tipo di allenatore è quello della riedizione per console corrente.

In seconda generazione la via è completa, documentata in forma di procedura e compatibile con ciò che l'utente possiede, ed è per questo la parte più utilizzabile della materia. Le guide della wiki dedicata coprono ogni versione e ogni lingua dei tre titoli, l'italiano compreso, e sono dichiarate compatibili sia con le cartucce originali sia con le riedizioni per console corrente \cite{timovm-ace-gen2}. La compatibilità con l'italiano non è un dettaglio: la catena di questo lavoro vincola ogni anello alla medesima lingua, e una via che funzionasse soltanto sull'inglese non servirebbe. L'allestimento si raggiunge presto nella partita, e la fonte dichiara fra trenta e sessanta minuti per installare un programma di cinquanta byte che permette di scrivere ed eseguire istruzioni arbitrarie senza dover più passare per le quantità degli oggetti o per i caratteri dei nomi delle scatole, che erano il vincolo dei metodi precedenti.

I codici già scritti che quella comunità pubblica sono la parte che riguarda la produzione \cite{glitchcity-mailwriter}. Tre di essi valgono da soli la lettura: uno consegna un esemplare mitico di livello cinque adattandone nome, identificativo e soprannome dell'allenatore perché il trasferimento lo accetti, con i valori individuali scelti da chi esegue; uno fa la stessa cosa per la linea di due iniziali della terza generazione, cioè per esemplari altrimenti non trasferibili da quella via; e uno riempie la scatola attiva chiedendo ogni volta specie e livello dell'esemplare successivo, ed è in sostanza un generatore interno al gioco. Accanto stanno codici che non producono esemplari ma cambiano le condizioni di una caccia, come rendere cromatico ogni incontro selvatico o imporgli un insieme di valori individuali, e un programma più grande con cui si leggono e si cambiano i valori in memoria, che è lo strumento con cui si verifica ciò che i codici hanno fatto invece di fidarsi.

I vincoli d'uso dichiarati vanno riportati perché sono il genere di dettaglio che fa fallire un tentativo: i codici che consegnano un esemplare pretendono una squadra di dimensione esatta, quello che riempie la scatola pretende la squadra piena e la scatola attiva non vuota, e dopo l'esecuzione i punti potenza delle mosse risultano azzerati e si sistemano depositando e ritirando l'esemplare.

In terza generazione la fonte non pubblica una procedura sola ma tre strumenti, cioè un generatore di codici da cui si compone il proprio, una raccolta di codici già scritti per tutte le lingue, e un editor di salvataggi pensato per chi lavora con i difetti sfruttabili \cite{ace-archive}. È una differenza di maturità e non di natura: in seconda generazione si eseguono codici scritti da altri, in terza si compone il proprio. La conseguenza per questo lavoro ha una forma paradossale che vale enunciare, perché contraddice l'attesa: la terza generazione, che è quella in cui il progetto sa già comporre il dato fuori dal gioco meglio che altrove, è anche quella in cui la via interna richiede più lavoro di allestimento.

\section{La manipolazione del generatore, che non è un difetto sfruttato}

Accanto alle vie precedenti sta una via che non impiega alcun difetto: manipolare il generatore pseudocasuale del gioco per far uscire l'esemplare che si vuole, con l'aspetto e i valori che si vogliono, senza modificare nulla. È l'alternativa legittima nel senso più stretto del termine, perché l'esemplare che ne esce è catturato per caso, dove il caso è stato scelto invece di essere atteso.

Va dichiarato un limite di questo lavoro e non delle fonti: l'estrattore che ha ridotto a testo le due pagine principali su questa tecnica ha conservato la navigazione e perso il contenuto, quindi di quelle voci oggi si sa che esistono e che cosa coprono, non che cosa insegnano. Il grezzo sta accanto al derivato e una seconda passata lo recupera senza scaricare di nuovo. La ragione per cui vale farla è che, una volta deciso di produrre tutto ciò che la scadenza rende irraggiungibile, la domanda su questa via non è più se sia lecita ma quanto costi in tempo, che è l'unica grandezza con cui compete.

Due voci laterali vanno registrate perché documentano punti d'ingresso che altrimenti si riscoprono. La prima è l'esecuzione di codice attraverso un oggetto contenitore in seconda generazione, che è la via storica precedente e che compare qui perché è il modo con cui si ottiene un oggetto malformato impiegato altrove. La seconda è lo sfruttamento del meccanismo di scambio fra prima e seconda generazione, che permette di portare in prima generazione esemplari di seconda mascherandoli, con la conseguenza che diventano esemplari malformati secondo una tabella di conversione della ROM: non serve alla collezione e va nella categoria delle curiosità, ma documenta un secondo punto d'ingresso.

\section{Che cosa ne discende per un lotto già prodotto}

La conseguenza pratica maggiore riguarda il lotto delle prime due generazioni, che oggi esiste come centosessantacinque file destinati a essere caricati in un salvataggio. Se la via interna funziona come la fonte dichiara, per una parte di quelle voci esiste un modo di ottenerle dentro il gioco invece di scriverle accanto, e la differenza è quella fra creare e ottenere. Non per tutte: i codici pubblicati coprono l'esemplare mitico, due linee di iniziali e il riempimento generico della scatola, mentre le voci con allenatore e soprannome storici richiederebbero codice proprio, che il programma di scrittura consente e che nessuno ha già scritto.

Restano tre verifiche, e l'ordine in cui conviene farle è quello del loro effetto. La prima è se un esemplare consegnato dal gioco per effetto di un codice sia trattato come ottenuto dal servizio di deposito, che è la domanda che decide se la via valga il suo costo. La seconda è l'aggiunta della voce mancante di prima generazione all'asse delle distribuzioni, con i suoi valori massimi e il suo identificativo. La terza è il recupero del contenuto delle pagine sulla manipolazione, oggi catalogate e non lette per un difetto della nostra estrazione.

\section{La via del nome di dominio, e un servizio che qualcuno ha rianimato}

La seconda via che non compone nulla è di natura completamente diversa e riguarda la quarta e la quinta generazione. I servizi in rete di quei giochi, chiusi nel 2014, sono stati ricostruiti da terzi e sono raggiungibili cambiando un solo parametro sulla console, cioè l'indirizzo del server dei nomi \cite{shacknews-dns}. Ne discendono due capacità distinte che vanno tenute separate perché hanno peso diverso: la consegna dei doni segreti, con un elenco lungo di esemplari che un utente dichiara di avere ricevuto su cinque giochi, e il sistema di scambio globale, dichiarato pienamente funzionante per la quarta generazione.

La differenza rispetto a tutto ciò che questo lavoro produce è la stessa enunciata all'inizio del capitolo, ed è qui nella sua forma più pura: il gioco riceve una carta vera attraverso il proprio canale di consegna, la riscatta con il proprio codice, e l'esemplare che ne esce è generato dal gioco. La correlazione fra valore di personalità e valori individuali non va imitata perché non è imitata, e nessun campo è scritto da noi. La cosa da provare cambia di conseguenza: non è più se un file somigli a un esemplare, è se un esemplare vero sia stato ottenuto per una via che il servizio non prevedeva.

Sullo stato attuale il lavoro distingue con cura ciò che ha verificato da ciò che ha soltanto letto, perché due delle tre fonti sono vecchie e in questo dominio sei anni sono molti. È verificato che il servizio sia vivo: la sua pagina, letta nel momento in cui la si scarica, dichiara contatori di esercizio aggiornati, cioè esemplari offerti e video di lotte registrati, che sono dati di funzionamento e non testo redazionale, e l'indirizzo che pubblica coincide con uno di quelli indicati dalle fonti più vecchie \cite{pkmnclassic}. Non è verificato che quel servizio distribuisca oggi i doni segreti: la sua pagina nomina lo scambio, i video e servizi generici, e non li nomina, mentre nelle fonti più vecchie la consegna dei doni viene dai progetti di emulazione generale della connessione e non da questo servizio, che vi si appoggia. Sono due funzioni distinte servite da due strati distinti, e concluderne che chi fornisce l'una fornisca anche l'altra sarebbe esattamente l'inferenza che questo lavoro marca e non promuove. Non è verificato, e importa ancora di più, quale catalogo la via distribuisca oggi: l'elenco disponibile è la testimonianza di un utente nel 2020 e non il catalogo del servizio.

Gli ostacoli pratici non sono piccoli e sono di tre nature. Quello tecnico riguarda la sola quarta generazione ed è l'anzianità della sua radio, che parla soltanto con reti prive di cifratura o con la più vecchia, oggi non più offerte dai punti di accesso: la via praticata dalle fonti è un punto di accesso mobile aperto, cioè una condizione da allestire con cura e per il tempo strettamente necessario. Quello di provenienza è la data, perché una carta ricevuta oggi scrive nell'esemplare la data di oggi, e per un evento chiuso nel 2013 quella data è precisamente la firma con cui la comunità riconosce questa via; per la quarta generazione il discorso è più favorevole, perché quel formato non porta marchio di origine e il deposito ne mostra la data del trasferimento, quindi una data recente non la distingue da un trasferimento tardivo legittimo. Quello di perimetro non è tecnico e la decisione non appartiene a chi scrive: percorrere la via significa collegare una console a un servizio non ufficiale che si sostituisce a uno dismesso, e sta accanto alla decisione sulla produzione senza coincidervi, perché là si trattava di scrivere byte e qui di usare un canale ricostruito.

\section{Le due decisioni che questa via rimetterebbe in discussione}

La prima riguarda i ventotto esemplari coreani della quarta generazione, il cui rifiuto era stato attribuito a un vincolo storico, cioè che quei giochi non si scambiassero direttamente con quelli internazionali e che il solo ponte fosse il sistema di scambio globale, allora chiuso. Se quel sistema è oggi vivo per la quarta generazione, la premessa su cui poggiava la scelta di allestire una stazione coreana in emulazione cambia. Due condizioni restano da verificare prima di poterlo affermare: che il sistema funzioni davvero fra un gioco coreano e uno internazionale e non soltanto fra due internazionali, e che accetti quelle voci, dato che il ponte storico rifiutava le uova e gli esemplari con un certo fiocco.

La seconda riguarda l'ordine di produzione, e la sua formulazione è la parte che conta. Se una parte dei doni di quarta e quinta generazione si può ricevere invece che comporre, quella parte non va composta, e non per risparmio di lavoro: un esemplare ricevuto dal gioco è migliore di uno composto su ogni dimensione che a questo lavoro interessi. La conseguenza operativa è che prima di continuare a comporre conviene sapere che cosa il canale distribuisca, il che è una misura e non una lettura.

I passi che ne discendono sono quattro e nessuno impegna il lavoro alla via. Il primo è misurare che cosa il canale offra oggi, leggendo il deposito pubblico del suo codice e il suo canale di conversazione, che sono raggiungibili con gli strumenti già in uso, e la domanda va posta stretta: se distribuisca doni oltre allo scambio, e se esista un catalogo di ciò che distribuisce. Il secondo, se il primo è positivo, è confrontare quel catalogo con le nostre due enumerazioni, che è lavoro da programma. Il terzo è una prova su un solo esemplare, con la quinta generazione perché non richiede il punto di accesso aperto, e con l'esame del risultato nel verificatore prima di qualunque trasferimento. Il quarto riguarda i coreani e chiede se il ponte ricostruito attraversi la barriera di lingua, che è una domanda a cui il codice del servizio può rispondere, perché quella barriera è una regola del gioco e non del server, e quindi il server la eredita o la ignora a seconda di come è fatto.

I primi due passi servono a sapere se valga la pena porsi la domanda di perimetro, che è la sola che non si risolve misurando.

Il primo dei quattro passi ha avuto una risposta parziale nel giro di lettura successivo, e va registrata con il suo grado di fiducia perché cambia l'ordine del lavoro senza autorizzare nulla. Una testimonianza del 2020, aggiornata da oltre mille commenti e confermata da due voci indipendenti del settembre 2026, elenca titolo per titolo i doni che il suo autore dichiara di avere ricevuto collegando la console al servizio ricostruito attraverso una diversa configurazione dei server dei nomi \cite{dns-eventi-gen45}. La risposta al primo passo è quindi affermativa su entrambe le sue metà: il canale distribuisce doni e non soltanto scambi, e un catalogo di ciò che distribuisce esiste, sia pure in forma di testimonianza invece che di elenco pubblicato.

Due dettagli di quell'elenco meritano di essere isolati perché toccano due voci che questo lavoro aveva classificato come irraggiungibili. Il primo è che il canale serve anche due doni che non furono mai distribuiti ufficialmente, e uno di essi è precisamente l'oggetto la cui assenza rendeva non conforme un esemplare che questo lavoro aveva registrato come non producibile: riceverlo dal canale non lo rende conforme, perché il verificatore non ha alcuna voce che gli corrisponda, ma sposta la questione da tecnica a di perimetro. Il secondo è che le voci mancate dai due testimoni sono le stesse in entrambe le testimonianze, il che suggerisce che l'insieme distribuito sia stabile e non casuale, e quindi misurabile.

Resta ciò che una testimonianza non può dare, e va detto perché è la differenza fra una misura e una speranza: che qualcuno abbia ricevuto quelle voci non dimostra che il canale le distribuisca a chiunque oggi, e nessuna delle due testimonianze dice nulla sui campi dell'esemplare ricevuto, che sono l'unica cosa che decide se il verificatore lo accetti. Il terzo dei quattro passi, cioè la prova su un solo esemplare esaminato prima di qualunque trasferimento, resta quindi necessario esattamente come prima.
